\scanpage{48}
quadratique non dégénérée sur $V_{\mathbb C}$,
\unsure{pour|par} laquelle $V_{\mathbb C}$ se décompose en
sous espaces isotropes
\[
(16)\qquad V_{\mathbb C}=\bigoplus_{\omega\in\underline\omega}V_{\mathbb C}^{\omega}
\]
qui sont aussi les sous-espaces propres
pour le groupe \inserted{$O^+(V)$} des rotations de $V$.
La donnée d’une orientation \inserted{$\omega$} de $V$
permet de définir un autom \unsure{d}
\[
(17)\qquad\left\{\begin{aligned}
&u_\omega\in\deleted{\unsure{\operatorname{Aut}}_{\mathbb R}(V)}\ O^+(V)\\
&u_\omega^2=-1
\end{aligned}\right.
\]
\margin{rotation d’un quart de tour dans le sens de l’orientation $\omega$}
d’où une structure complexe sur
$V$\inserted{(qui devient un \unsure{vect.} complexe de rang $1$, soit $V_{(\omega)}$)}, en faisant opérer $i\in\mathbb C$ par
\[
i\mathbin{\underset{\omega}{\cdot}}x=u_\omega x
\]
Mais par la définition universelle du
complexifié, la \deleted{structure} application
$\mathbb R$-linéaire $V\longrightarrow V_\omega$ se prolonge
canoniquement en une application

\margin{47}
